\documentclass[10pt,a4paper]{amsart}
\usepackage[margin=19mm]{geometry}
\usepackage{amsmath,amssymb,mathtools}
\usepackage[scr=rsfs,cal=euler]{mathalfa}
\usepackage{xfrac}
\usepackage[unicode,hidelinks,bookmarks=false]{hyperref}
\newcommand{\ie}{i.e.\ }
\newcommand{\cf}{cf.\ }
\newcommand{\resp}{respectively\ }
\newcommand{\Subsection}{\subsection}
\providecommand{\nobreakdash}{\nobreak\hbox{-}}
\setlength{\emergencystretch}{2em}
\multlinegap0pt
\usepackage{mathtools}
\newtheorem{corollary}{Corollary}
\title{On finite groups with exactly one noncommutator}
\author{Saveliy V. Skresanov}
\date{Source: arXiv:2509.17587v1 (CC BY 4.0)}
\begin{document}
\maketitle

\noindent\emph{Source and licence.} On finite groups with exactly one noncommutator. Version arXiv:2509.17587v1 (CC BY 4.0), distributed by arXiv under the Creative Commons Attribution 4.0 International licence, http://creativecommons.org/licenses/by/4.0/, as recorded in the arXiv licence metadata for this exact version and shown on the abstract page https://arxiv.org/abs/2509.17587v1. Full licence notice: \texttt{licenses/arxiv-2509.17587-CC-BY-4.0.txt}.\par\smallskip

\noindent\textit{Editorial scope.} Counted unit: the article's corollary corollary (source0.tex lines 72--74). The article's complete original proof of it is delivered with it (source0.tex lines 75--95, one \texttt{proof} environment of 21 lines). No mathematics is written, repaired or re-derived: the statement and the proof are the article's own lines, in the article's own notation, and no retained line is substituted or re-flowed.  No further source-local context is needed: every cross-reference of every retained line resolves inside the two delivered spans.  Nothing was omitted from any retained span, so the package carries no omission record.  No retained line contains a citation, so the wrapper carries no bibliography and no bibliography entry.  No retained line contains a figure environment, an image-inclusion command or drawing code, so no image is delivered and the package declares no asset dependency.  Editorial formatting only, in addition: an AMS \texttt{amsart} preamble that restates the article's own declarations and macros used by the retained lines, namely the environments \texttt{corollary} on separate counters, together with the standard packages those lines need.  The retained environments print in this document as Corollary 1 in order of appearance, whereas the article prints the same items with the numbering of its own full document.  The complete native source of the article as distributed in the arXiv e-print for this exact version is bundled unchanged as \texttt{sources/185467/source0.tex}.
\begin{corollary}
	There are infinitely many finite groups with exactly one noncommutator.
\end{corollary}

\begin{proof}
	Let \( G \) be a finite nonabelian group with exactly one noncommutator~\( t \).
	Let \( T \) be a normal subgroup of \( G \times G \) generated by \( (t, t) \), and set \( K = (G \times G)/T \).
	Notice that \( |K| > |G| \). We will prove that \( K \) also has exactly one noncommutator.

	Suppose that \( (t, 1)T \in K \) is a commutator. Then there exist elements \( a, b \in G \times G \)
	such that \( (t, 1) \in [a, b]T \). If \( a = (a_1, a_2) \) and \( b = (b_1, b_2) \), then
	\( (t, 1) = ([a_1, b_1], [a_2, b_2])T \). Hence either \( t = [a_1, b_1] \), which is impossible since \( t \)
	is not a commutator in \( G \), or \( 1 = [a_2, b_2]t \). But that implies \( t = [a_2, b_2] \) which is again a contradiction.
	Therefore \( (t, 1) \) is not a commutator in \( K \).

	Let \( (x, y)T \in K \) be an element distinct from \( (t, 1)T \).
	We may assume that \( x, y \neq t \). Indeed, if \( x = t \), then \( y \neq 1 \), since \( (x, y) \not\in (t, 1)T \).
	So we may replace the representative \( (x, y) \) by \( (xt, yt) = (1, yt) \), where \( 1, yt \neq t \).
	Similar reasoning applies in the case when \( y = t \). Now, if \( x, y \neq t \), there exist elements \( a_1, a_2, b_1, b_2 \in G \)
	such that \( x = [a_1, b_1] \), \( y = [a_2, b_2] \), and hence \( (x, y)T = [a, b]T \) is a commutator in \( K \),
	where \( a = (a_1, a_2) \), \( b = (b_1, b_2) \).

	Given a nonabelian group \( G \) with one noncommutator we constructed a larger group \( K \) with one noncommutator.
	By applying this procedure repeatedly to the original group \( G_{MacHale} \), we obtain an infinite series of groups with one noncommutator.
\end{proof}

\end{document}
